\documentclass[11pt,a4paper]{article}
\usepackage[utf8]{inputenc}
\usepackage{amsmath,amsfonts,amssymb}
\usepackage{physics}
\usepackage{geometry}
\usepackage{listings}
\usepackage{xcolor}
\usepackage{hyperref}
\usepackage{float}
\usepackage{graphicx}

\geometry{margin=1in}

% Code formatting
\lstset{
    basicstyle=\ttfamily\footnotesize,
    commentstyle=\color{green!60!black},
    keywordstyle=\color{blue},
    stringstyle=\color{red},
    showstringspaces=false,
    breaklines=true,
    frame=single,
    numbers=left,
    numberstyle=\tiny\color{gray}
}

\title{Twisted Angle Planar Wave (TAPW) Method Implementation\\
Band Structure Calculation Procedure}
\author{GRABNES contributors}
\date{Pre-release technical note}

\begin{document}

\maketitle

\section{Introduction}

The Twisted Angle Planar Wave (TAPW) method is implemented in the \texttt{DiagBands} routine for calculating electronic band structures of moiré systems. This document provides a comprehensive overview of the TAPW procedure from entry into \texttt{DiagBands} through generation of the \texttt{generate.bands} file, excluding Chern number calculations.

\section{Overview of TAPW Flow}

The TAPW calculation follows this sequence:
\begin{enumerate}
    \item \textbf{Parameter Setup}: Read TAPW parameters and initialize coordinate systems
    \item \textbf{G-vector Generation}: Create reciprocal space grid around reference K-point
    \item \textbf{Basis Construction}: Build TAPW basis functions using plane waves
    \item \textbf{Hamiltonian Projection}: Project tight-binding Hamiltonian onto TAPW basis
    \item \textbf{Diagonalization}: Solve eigenvalue problem for each k-point
    \item \textbf{Output Generation}: Write band structure to \texttt{generate.bands}
\end{enumerate}

\section{Detailed Implementation}

\subsection{Entry Point and Parameter Setup}

The TAPW calculation begins when \texttt{useTAPW = .true.} in the \texttt{DiagBands} routine:

\begin{lstlisting}[language=Fortran,caption=TAPW Entry Point]
! In DiagBands routine
else if (useTAPW) then
   call MIO_Print('calling TAPW routine for '//trim(num2str(ip)),'diag')
   call DiagH0TAPW(nAt, nspin, is, ELoc, KptsLoc, ucell, H0, 
                   maxNeigh, hopp, NList, Nneigh, neighCell, neig)
end if
E(:, is, ip) = ELoc
\end{lstlisting}

Key parameters are read from input:
\begin{itemize}
    \item \texttt{Diag.N\_G}: Number of G-vectors (\texttt{tapwNG})
    \item \texttt{Diag.MoireAngle}: Rotation angle for graphene lattice
\end{itemize}

\subsection{Coordinate System Construction}

\subsubsection{Graphene Reciprocal Lattice}

The code constructs rotated graphene reciprocal vectors \texttt{rG}:

\begin{equation}
\mathbf{a}_1 = a_G \begin{pmatrix} 1 \\ 0 \end{pmatrix}, \quad
\mathbf{a}_2 = a_G \begin{pmatrix} \cos(60°) \\ \sin(60°) \end{pmatrix}
\end{equation}

After rotation by \texttt{moireAngle}:
\begin{equation}
\mathbf{a}'_i = \begin{pmatrix} \cos\theta & -\sin\theta \\ \sin\theta & \cos\theta \end{pmatrix} \mathbf{a}_i
\end{equation}

Reciprocal vectors are computed as:
\begin{equation}
\mathbf{rG} = 2\pi \cdot (\mathbf{A}^T)^{-1}
\end{equation}

\begin{lstlisting}[language=Fortran,caption=Reciprocal Lattice Construction]
! Build graphene lattice vectors (60-degree structure)
sGlattice(1,1) = 1.0_dp * aG                    ! (aG, 0)
sGlattice(2,1) = 0.0_dp
sGlattice(1,2) = cos(60.0_dp*pi/180.0_dp) * aG  ! (aG*cos60°, aG*sin60°)
sGlattice(2,2) = sin(60.0_dp*pi/180.0_dp) * aG

! Apply rotation
rotation_angle = moireAngle * pi / 180.0_dp
sGlattice = matmul(reshape((/ cos_rot, sin_rot, -sin_rot, cos_rot /), 
                           (/2,2/)), sGlattice)

! Convert to reciprocal space: rG = 2π * inverse(lattice.T)
rG = 2.0_dp * pi * transpose(sGlattice_inv)
\end{lstlisting}

\subsubsection{Reference K-point}

The K-point in graphene's Brillouin zone is calculated as:
\begin{equation}
\mathbf{k}_{ref} = \frac{2}{3}\mathbf{b}_1 + \frac{1}{3}\mathbf{b}_2
\end{equation}

\begin{lstlisting}[language=Fortran,caption=K-point Calculation]
real(dp), parameter :: refF(2) = (/ 2.0_dp/3.0_dp, 1.0_dp/3.0_dp /)
k_ref = refF(1)*rG(:,1) + refF(2)*rG(:,2)
\end{lstlisting}

\subsection{G-vector Generation}

The TAPW method generates a set of reciprocal lattice vectors around the reference K-point. Two methods are implemented:

\subsubsection{Hexagonal Shell Method}

\begin{lstlisting}[language=Fortran,caption=G-vector Generation]
call generate_shifted_G_list(rcell, k_ref, NGrange, Gx, Gy, NG)
\end{lstlisting}

This generates G-vectors as:
\begin{equation}
\mathbf{G}_{n_1,n_2} = n_1 \mathbf{b}_{moire,1} + n_2 \mathbf{b}_{moire,2}
\end{equation}

The vectors are:
\begin{enumerate}
    \item Generated on a grid around \texttt{k\_ref}
    \item Sorted by distance from the nearest lattice point
    \item Selected to keep the \texttt{N\_G} closest vectors
    \item Filtered for uniqueness
\end{enumerate}

\subsubsection{Triangular Truncation Method}

Alternative method following Python implementation:
\begin{equation}
\mathbf{G}_{frac} = \begin{pmatrix} j + k_{ref,x} \\ i + k_{ref,y} \end{pmatrix}
\end{equation}

Selected within triangular Brillouin zone boundaries.

\subsection{TAPW Basis Construction}

\subsubsection{Basis Functions}

The TAPW basis functions are plane waves modulated by atomic orbitals:
\begin{equation}
\phi^{\alpha}_{\mathbf{G}}(\mathbf{r}) = \sum_{\mathbf{R}} e^{i(\mathbf{k} + \mathbf{G}) \cdot (\mathbf{R} + \boldsymbol{\tau}_\alpha)} \psi_\alpha(\mathbf{r} - \mathbf{R} - \boldsymbol{\tau}_\alpha)
\end{equation}

where:
\begin{itemize}
    \item $\alpha$ labels atomic species/sublattices
    \item $\mathbf{G}$ are the selected reciprocal lattice vectors
    \item $\mathbf{R}$ are lattice translation vectors
    \item $\boldsymbol{\tau}_\alpha$ are atomic positions within unit cell
\end{itemize}

\subsubsection{X-Matrix Construction}

The transformation matrix $X$ connects atomic orbitals to TAPW basis:

\begin{lstlisting}[language=Fortran,caption=X-Matrix Construction]
call build_X(X, x, y, label, Gx, Gy, N, NG, Nlabel)
\end{lstlisting}

Matrix elements:
\begin{equation}
X_{i,\alpha\mathbf{G}} = e^{i\mathbf{G} \cdot \mathbf{r}_i} \delta_{\alpha,\text{species}(i)}
\end{equation}

Implemented as:
\begin{lstlisting}[language=Fortran,caption=X-Matrix Elements]
do i = 1, N
   do iG = 1, NG
      do iLabel = 1, Nlabel
         if (label(i) == iLabel) then
            idx = (iLabel-1)*NG + iG
            phase = Gx(iG)*x(i) + Gy(iG)*y(i)
            X(i, idx) = cmplx(cos(phase), sin(phase), dp)
         end if
      end do
   end do
end do
\end{lstlisting}

\subsection{Hamiltonian Projection}

\subsubsection{Sparse Matrix Construction}

The tight-binding Hamiltonian is first constructed in sparse CSR format:

\begin{lstlisting}[language=Fortran,caption=Sparse Hamiltonian]
call initialize_sparse_matrix(N, maxN, H0, hopp, NList, Nneigh, 
                             neighCell, ns, is, KLoc, cell, 
                             row_ptr, col_ind, values, sigma)
\end{lstlisting}

Matrix elements include:
\begin{equation}
H_{ij} = \epsilon_i \delta_{ij} + \sum_{\mathbf{R}} t_{ij}(\mathbf{R}) e^{i\mathbf{k} \cdot \mathbf{R}}
\end{equation}

where $t_{ij}(\mathbf{R})$ are hopping parameters and $\epsilon_i$ are on-site energies.

\subsubsection{TAPW Projection}

The projected Hamiltonian in TAPW space:
\begin{equation}
H^{TAPW} = X^\dagger H X
\end{equation}

\begin{lstlisting}[language=Fortran,caption=Hamiltonian Projection]
call transform_sparse_hamiltonian(N, M, row_ptr, col_ind, values, X, Hproj)
\end{lstlisting}

Matrix dimensions: $M = N_G \times N_{label}$ where $N_{label}$ typically equals 2 (sublattices A and B).

\subsection{Eigenvalue Solution}

The projected Hamiltonian is diagonalized using LAPACK:

\begin{equation}
H^{TAPW} \psi^{TAPW}_n = E_n \psi^{TAPW}_n
\end{equation}

\begin{lstlisting}[language=Fortran,caption=Diagonalization]
allocate(eigvals(M))
lwork = 2*M
allocate(ZWorkLoc(2*M), DWorkLoc(3*M))

call ZHEEV('V','U', M, Hproj, M, eigvals, ZWorkLoc, lwork, DWorkLoc, info)

if (info /= 0) then
   print *, "Diagonalization failed: ZHEEV info =", info
   stop
end if
\end{lstlisting}

\subsection{K-point Loop and Progress Tracking}

The TAPW calculation is performed for each k-point in the path:

\begin{lstlisting}[language=Fortran,caption=K-point Loop]
do ip = 1, ptsTot
   ! Progress tracking for TAPW (show every 10%)
   if (useTAPW .and. modulo(ip, max(1, int(ptsTot/10))) == 0) then
      call MIO_Print('TAPW progress: '//trim(num2str(nint(100.0_dp*ip/ptsTot)))//
                     '% ('//trim(num2str(ip))//'/'//trim(num2str(ptsTot))//
                     ' k-points)', 'diag')
   end if
   
   do is = 1, nspin
      KptsLoc = Kpts(:, ip)
      call DiagH0TAPW(nAt, nspin, is, ELoc, KptsLoc, ucell, H0, 
                      maxNeigh, hopp, NList, Nneigh, neighCell, neig)
      E(:, is, ip) = ELoc
   end do
end do
\end{lstlisting}

\section{Output Generation}

\subsection{Band Structure File}

After all k-points are calculated, the results are written to \texttt{generate.bands}:

\begin{lstlisting}[language=Fortran,caption=Band Output]
! Write TAPW matrix size header
if (useTAPW) then
   write(uu,'(A,I0)') '# TAPW matrix size: M = ', M_tapw
end if

! Write eigenvalues for each k-point
do ip = 1, ptsTot
   do is = 1, nspin
      do i = 1, nAt
         write(uu, *) E(i, is, ip)
      end do
   end do
end do
\end{lstlisting}

\subsection{Debug Output Files}

Several debug files are generated:
\begin{itemize}
    \item \texttt{brillouin\_zones\_debug.dat}: Brillouin zone vertices and K-points
    \item \texttt{moire\_bz\_debug.dat}: Moiré BZ information
    \item \texttt{g\_vectors\_debug.dat}: Generated G-vectors
    \item \texttt{kpath\_debug\_*}: K-path coordinates
\end{itemize}

\section{Key Equations Summary}

\subsection{Core TAPW Equations}

\begin{align}
\text{Basis functions:} \quad & \phi^{\alpha}_{\mathbf{G}}(\mathbf{r}) = \sum_{\mathbf{R}} e^{i(\mathbf{k} + \mathbf{G}) \cdot (\mathbf{R} + \boldsymbol{\tau}_\alpha)} \psi_\alpha(\mathbf{r} - \mathbf{R} - \boldsymbol{\tau}_\alpha) \\
\text{Transformation matrix:} \quad & X_{i,\alpha\mathbf{G}} = e^{i\mathbf{G} \cdot \mathbf{r}_i} \delta_{\alpha,\text{species}(i)} \\
\text{Projected Hamiltonian:} \quad & H^{TAPW} = X^\dagger H X \\
\text{Eigenvalue equation:} \quad & H^{TAPW} \psi^{TAPW}_n = E_n \psi^{TAPW}_n
\end{align}

\subsection{Coordinate Systems}

\begin{align}
\text{Graphene lattice:} \quad & \mathbf{a}_1 = a_G(1,0), \quad \mathbf{a}_2 = a_G(\cos 60°, \sin 60°) \\
\text{Reciprocal lattice:} \quad & \mathbf{rG} = 2\pi \cdot (\mathbf{A}^T)^{-1} \\
\text{K-point:} \quad & \mathbf{k}_{ref} = \frac{2}{3}\mathbf{b}_1 + \frac{1}{3}\mathbf{b}_2 \\
\text{G-vectors:} \quad & \mathbf{G}_{n_1,n_2} = n_1 \mathbf{b}_{1} + n_2 \mathbf{b}_{2}
\end{align}

\section{Implementation Parameters}

\begin{table}[H]
\centering
\begin{tabular}{|l|l|l|}
\hline
\textbf{Parameter} & \textbf{Variable} & \textbf{Description} \\
\hline
\texttt{Diag.N\_G} & \texttt{tapwNG} & Number of G-vectors (default: 180) \\
\texttt{Diag.MoireAngle} & \texttt{moireAngle} & Rotation angle in degrees \\
\texttt{refF} & \texttt{(2/3, 1/3)} & Fractional K-point coordinates \\
\texttt{Nlabel} & 2 & Number of sublattices (A, B) \\
\texttt{M} & \texttt{NG × Nlabel} & TAPW matrix dimension \\
\hline
\end{tabular}
\caption{Key TAPW Parameters}
\end{table}

\section{Computational Complexity}

\begin{itemize}
    \item \textbf{G-vector generation}: $O(N_G \log N_G)$ for sorting
    \item \textbf{X-matrix construction}: $O(N \times N_G \times N_{label})$
    \item \textbf{Hamiltonian projection}: $O(N_{sparse} \times M)$ where $N_{sparse}$ is number of non-zero elements
    \item \textbf{Diagonalization}: $O(M^3)$ per k-point
    \item \textbf{Total per k-point}: $O(M^3)$ dominated by diagonalization
\end{itemize}

For typical parameters ($N_G = 180$, $N_{label} = 2$), $M = 360$, making diagonalization the computational bottleneck.

\section{Verification and Debugging}

The implementation includes comprehensive verification:
\begin{itemize}
    \item Direct-reciprocal lattice duality checks
    \item G-vector uniqueness and ordering verification  
    \item Hamiltonian Hermiticity validation
    \item Matrix norm consistency checks
    \item Coordinate system orientation verification
\end{itemize}

\section{Conclusion}

The TAPW implementation provides a systematic approach to calculating electronic band structures in moiré systems by:
\begin{enumerate}
    \item Constructing appropriate reciprocal space grids
    \item Building plane wave basis functions
    \item Projecting tight-binding Hamiltonians onto reduced subspaces
    \item Efficiently solving eigenvalue problems
    \item Generating standard band structure output
\end{enumerate}

The method effectively reduces the computational complexity from the full supercell size to a manageable projected space while maintaining the essential physics of the moiré system.

\end{document}
